
\paragraph{\textbf{Remarks 6.11.}}
\SGARefTarget{sga3:xiv:6.11}
\textup{(a)} Rosenlicht gave examples of twisted forms of
\(\mathbb G_a\) over a nonperfect field having only finitely many
rational points; \emph{a fortiori}, they are not unirational.
Chevalley, on the other hand, gave an example of a torus over a
field of characteristic zero that is not rational. By contrast,
Chevalley's theory of semisimple groups shows that, over an
algebraically closed field, every smooth connected affine
algebraic group is rational. In any case, the question of
unirationality arises only for affine algebraic groups: by
Chevalley's structure theorem, a unirational algebraic group is
necessarily affine.

\textup{(b)} With the notation of \SGARefLink{sga3:xiv:6.6}{6.6}, it is tempting to seek a
criterion for the unirationality of \(G\) in terms of \(G/H\),
which is unipotent when \(G\) is affine. Clearly \(G/H\) must be
unirational; is that condition also sufficient? An example of
Rosenlicht shows that a smooth connected unipotent algebraic
group \(U\) may be rational without being \(k\)-unipotent.

\textup{(c)} Over a finite field \(k\), it would be interesting
to study questions of ``density'' such as the following one,
raised by Rosenlicht: if \(G\) is a smooth connected algebraic
group over \(k\), is \(G\) generated by its Cartan
subgroups?\footnote{This question was subsequently answered in
the affirmative by Steinberg.}

\begin{quote}\small
\textbf{Editors' note.}
The editors did not
identify this result in Steinberg's \emph{Collected Papers}; here
is a proof based on the Bruhat decomposition.

Let \(G/k\) be semisimple over the finite field \(k\), and let
\(G^\sharp\) be the subgroup generated by its \(k\)-subtori. The
question is stable under central isogeny and restriction of
scalars, so, as in \SGARefLink{sga3:xiv:7:lemma1}{Lemma~1} of the Appendix below, it is enough to
treat the geometrically simple case. The group \(G\) is
quasi-split by Lang's theorem and admits a Killing pair
\((T,B)\). By definition \(T\subset G^\sharp\), hence \(T\)
normalizes \(G^\sharp\). The big cell
\[
  R_u(B^-)\,T\,R_u(B)
\]
shows that it suffices to prove \(R_u(B)\subset G^\sharp\).
Let \(S\) be the maximal split \(k\)-torus of \(T\). For the
relative root system \(\Phi(G,S)\), choose a basis
\(\Delta_k\), and denote by \(U(\alpha)\) the unipotent subgroup
associated with \(\alpha\)
\textup{(A.~Borel, \emph{Linear Algebraic Groups}, second
edition, 1991, Proposition~21.9)}. Since \(R_u(B)\) is generated
by the \(U(\alpha)\), \(\alpha\in\Delta_k\), it is enough to
show \(U(\alpha)\subset G^\sharp\).

The classification supplies a quasi-split semisimple subgroup
\(G_\alpha\) of type \(A_1\), \({}^{2}A_1\), \({}^{3}A_1\), or
\({}^{2}A_2\), with
\[
  U(\alpha)\subset G_\alpha\subset G.
\]
Thus it is enough to consider \(G=\operatorname{PGL}_2\) or
\(G=\operatorname{SU}_3(K)\), where \(K/k\) is the unique
quadratic extension.

For \(\operatorname{PGL}_2\), since \(G^\sharp\) contains the
standard split torus \(T\), the possibilities up to
\(G(k)\)-conjugacy are
\[
  G^\sharp=T,\qquad G^\sharp=B,\qquad
  G^\sharp=\operatorname{PGL}_2.
\]
The first is excluded because \(G^\sharp\) contains the
\(k\)-torus \(R_{K/k}(\mathbb G_m)/\mathbb G_m\). The preceding
discussion also shows that \(G^\sharp=B\) forces
\(G^\sharp=\operatorname{PGL}_2\). Hence
\(G^\sharp=\operatorname{PGL}_2\).

For \(G=\operatorname{SU}_3(K)\), one has
\(\Delta_k=\{\alpha\}\), and the possibilities up to conjugacy
are
\[
 \begin{split}
  G^\sharp={}&T,\quad B,\quad U(2\alpha)\rtimes T,\\
             &\langle U(2\alpha),U(-2\alpha),T\rangle
               =\operatorname{SL}_2\cdot T,\quad\text{or }G.
 \end{split}
\]
The case \(B\) is excluded as for \(\operatorname{PGL}_2\). The
case \(\operatorname{SL}_2\cdot T\) is excluded because
\(G^\sharp\) contains the \(k\)-torus
\[
  R^1_{K/k}\!\left(R^1_{K_3/K}(\mathbb G_m)\right),
\]
where \(K_3/K\) is the degree-three extension, and the isogeny
class of this torus is irreducible. The same torus excludes
\(T\) and \(U(2\alpha)\rtimes T\). Hence \(G^\sharp=G\).
\end{quote}

For Rosenlicht's question one may reduce to the affine case by
dividing by the center. In the semisimple case, an affirmative
answer would also follow from a strengthening of Chevalley's
existence result for regular points mentioned in \SGARefLink{sga3:xiv:3.20}{3.20}: one would
need a regular element of \(\mathfrak g\) outside
\(\mathfrak h=\operatorname{Lie}(H)\), where \(H\ne G\) is a
prescribed smooth algebraic subgroup of an adjoint semisimple
group \(G\) over the finite field \(k\).
